%!TEX root = ../../../../exa-ma-d7.1.tex
\VThreeApplicationDraft{app-feelpp-hifimagnet-optcontrol}
{HiFiMagnet robust optimal operation and control}
{C. Prud'homme; V. Chabannes}
{WP1, WP2, WP3, WP5, WP6}{extended-mini-app}
{Determine coil-group currents and cooling controls that reach a target magnetic field while limiting power, temperature, voltage and actuator constraints.}
{Feel++, PETSc/TAO, HPDDM, Sage-HPC; MPI, Kokkos and HIP}{spec-only}
\subsection{Description of the benchmark}
Optimal-operation application: determine coil-group currents and cooling controls that reach a target magnetic field with minimum power and thermal stress while satisfying temperature, voltage, power, total-current and actuator constraints.
\paragraph{Specification boundary.}
The output is an operating/control vector and its engineering margins, not
calibrated material parameters.  \texttt{BM-FEEL-HIFI-CTRL-01} is a proposed
verification/scalability record based on Taylor-gradient tests, target-field
error, constraint violation and KKT residual.
\todo{Supply the reviewed objectives/constraints, uncertainty scenarios,
implementation commit, reference solution, machine and raw evidence.}

\subsection{Benchmarking methodology}


\begin{description}
  \item[Method-Algorithm WP1:] nonlinear thermo-electric finite elements, coupled PDE-scalar formulation, integral total-current boundary condition
  \item[Method-Algorithm WP2:] ROM, hyper-reduction, surrogate-assisted optimization, FOM/ROM switching
  \item[Method-Algorithm WP3:] nonlinear block solvers, repeated state/adjoint solves, PETSc/HPDDM preconditioning
  \item[Method-Algorithm WP5:] PDE-constrained optimal operation and control: target-field tracking, power and thermal objectives; adjoint gradients; bound/general constraints; multiple coil-current and cooling controls; deterministic and robust formulations
  \item[Method-Algorithm WP6:] uncertain material and cooling parameters, stochastic and robust optimization, expectation and CVaR objectives
  \item[Framework:] Feel++, PETSc/TAO, HPDDM, Sage-HPC
  \item[parallel\_framework:] MPI, Kokkos, GPU - HIP
  \item[Benchmark scope:] multi-node, multi-gpu, full-system, solver-scaling, ensemble-mode
  \item[metrics:] Taylor-gradient-test, target-field-error, electrical-power, peak-temperature, constraint-violation, KKT-residual, time-to-optimum, PDE-solves-to-solution, TAO-iterations, strong-scalability, weak-scalability, scenario-throughput
\end{description}
\todo[app-feelpp-hifimagnet-optcontrol / benchmarking methodology]{Benchmarking methodology to be supplied by contributors.}

\subsection{Input/Output Dataset Description}

\subsubsection{Input Data}
Gmsh, JSON config, nonlinear material laws, target field, current/cooling bounds, uncertainty scenarios

\subsubsection{Output Data}
Optimal control vector z* (coil-group currents and cooling controls), optimal state/adjoint fields, magnetic field, power and peak-temperature predictions, voltage and constraint margins, KKT and objective histories, Pareto/robustness reports

\subsection{Results summary}
\todo[app-feelpp-hifimagnet-optcontrol / results summary]{Results summary to be supplied by contributors.}
